% Part: sets-functions-relations
% Chapter: sets
% Section: basics
% Hindi translation (hi-Deva-IN) of Open Logic commit 9620cc73f9c8e0ad003c514a5d3748f29611c4c0.

\documentclass[../../../include/open-logic-section]{subfiles}

\begin{document}

\olfileid[hi]{sfr}{set}{bas}
\olsection{अवयव-निर्धारित समानता}

\emph{समुच्चय} वस्तुओं का ऐसा संग्रह है जिसे एक ही वस्तु माना जाता है।
समुच्चय बनाने वाली वस्तुओं को उसके \emph{अवयव} या \emph{सदस्य} कहते हैं।
यदि $x$ किसी समुच्चय~$A$ का !!a{element} है, तो हम $x \in A$ लिखते हैं;
यदि नहीं, तो $x \notin A$ लिखते हैं। जिस समुच्चय का कोई !!{element} नहीं है,
उसे \emph{रिक्त समुच्चय} कहते हैं और ``$\emptyset$'' से निरूपित करते हैं।

\begin{explain}
समुच्चय को हम \emph{कैसे निर्दिष्ट} करते हैं, उसके !!{element} किस
\emph{क्रम} में रखते हैं, अथवा उसके !!{element} को \emph{कितनी बार}
गिनते हैं---इनमें से कोई बात महत्त्वपूर्ण नहीं है। केवल यह महत्त्वपूर्ण है कि
उसके !!{element} कौन-से हैं। इस विचार को हम निम्न सिद्धांत में सटीक रूप देते हैं।
\end{explain}

\begin{defn}[विस्तारिता]
  यदि $A$ और $B$ समुच्चय हैं, तो $A = B$ तभी और केवल तभी जब
  $A$ का प्रत्येक !!{element}, $B$ का भी !!a{element} हो, और इसके विपरीत।
\end{defn}

विस्तारिता कुछ संकेतनों का औचित्य सिद्ध करती है। सामान्यतः, जब हमारे पास
$a_{1}$, \dots, $a_{n}$ वस्तुएँ हों, तब $\{a_{1}, \dots, a_{n}\}$
\emph{वह} समुच्चय है जिसके !!{element} $a_1, \ldots, a_n$ हैं। हम
``\emph{वह}'' शब्द पर बल देते हैं, क्योंकि विस्तारिता बताती है कि ऐसा केवल
\emph{एक} समुच्चय हो सकता है। वास्तव में, विस्तारिता निम्न समानता का भी
औचित्य सिद्ध करती है:
  \[
    \{a, a, b\} = \{a, b\} = \{b,a\}.
  \]
इससे यही बात पुष्ट होती है कि समुच्चयों पर विचार करते समय हमें उनके
!!{element} के क्रम या किसी !!{element} को कितनी बार लिखा गया है, इसकी
परवाह नहीं होती।

\begin{tagblock}{novice}
\begin{ex}
जब भी आपके पास कुछ वस्तुएँ हों, आप उन्हें एक समुच्चय में एकत्र कर सकते हैं।
उदाहरण के लिए, रिचर्ड के भाई-बहनों का समुच्चय एक व्यक्ति को समाहित करता है;
हम इसे $S=\{\textrm{रुथ}\}$ लिख सकते हैं। $4$ से छोटे धनात्मक पूर्णांकों का
समुच्चय $\{1, 2, 3\}$ है, किंतु इसे $\{3, 2, 1\}$ या यहाँ तक कि
$\{1, 2, 1, 2, 3\}$ भी लिखा जा सकता है। विस्तारिता के अनुसार ये सभी एक ही
समुच्चय हैं। $\{1, 2, 3\}$ का प्रत्येक !!{element}, $\{3, 2, 1\}$ का भी
!!a{element} है (और $\{1, 2, 1, 2, 3\}$ का भी), तथा इसके विपरीत।
\end{ex}
\end{tagblock}

अक्सर हम किसी समुच्चय को उसके !!{element} के किसी साझा गुणधर्म द्वारा
निर्दिष्ट करेंगे। इसके लिए हम निम्न संक्षिप्त संकेतन का उपयोग करेंगे:
$\Setabs{x}{\phi(x)}$, जहाँ $\phi(x)$ उस गुणधर्म को दर्शाता है जो $x$ में
होना चाहिए ताकि उसे समुच्चय के !!{element} में गिना जाए।

\begin{tagblock}{novice}
\begin{ex}
अपने उदाहरण में हम $S$ को इस प्रकार भी निर्दिष्ट कर सकते थे:
\[
S = \Setabs{x}{x \text{ रिचर्ड का भाई या बहन है}}.
\]
\end{ex}
\end{tagblock}

\begin{tagblock}{math}
\begin{ex}
किसी संख्या को \emph{परिपूर्ण} तभी और केवल तभी कहते हैं जब वह अपने उचित
भाजकों के योग के बराबर हो (अर्थात् ऐसी संख्याएँ जो उसे पूरा-पूरा विभाजित करती
हैं, पर स्वयं उस संख्या के समान नहीं हैं)। उदाहरणतः, $6$ परिपूर्ण है क्योंकि
उसके उचित भाजक $1$, $2$ और~$3$ हैं तथा $6 = 1 + 2 + 3$। वास्तव में,
$10$ से छोटा एकमात्र धनात्मक पूर्णांक जो परिपूर्ण है, $6$ है। अतः विस्तारिता
का उपयोग करके हम कह सकते हैं:
\[
	\{6\} = \Setabs{x}{x\text{ परिपूर्ण है और }0 \leq x \leq 10}
\]
दाईं ओर के संकेतन को हम ``उन $x$ का समुच्चय जिनके लिए $x$ परिपूर्ण है और
$0 \leq x \leq 10$'' पढ़ते हैं। यहाँ की समानता इस बात की पुष्टि करती है कि
समुच्चयों पर विचार करते समय हमें इस बात की परवाह नहीं होती कि वे कैसे निर्दिष्ट
किए गए हैं। अधिक सामान्यतः, विस्तारिता सुनिश्चित करती है कि ऐसे $x$ का, जिनके
लिए $\phi(x)$ सत्य है, सदैव केवल एक ही समुच्चय होता है। इसलिए विस्तारिता
$\Setabs{x}{\phi(x)}$ को उन $x$ का \emph{वह} समुच्चय कहने का औचित्य
सिद्ध करती है जिनके लिए~$\phi(x)$ सत्य है।
\end{ex}
\end{tagblock}

विस्तारिता हमें यह दिखाने की विधि देती है कि दो समुच्चय समान हैं: $A = B$
दिखाने के लिए यह दिखाइए कि जब भी $x \in A$ हो, तब $x \in B$ भी हो, और
जब भी $y \in B$ हो, तब $y \in A$ भी हो।

\begin{prob}
सिद्ध कीजिए कि अधिकतम एक ही रिक्त समुच्चय हो सकता है; अर्थात् दिखाइए कि यदि
$A$ और $B$ ऐसे समुच्चय हैं जिनका कोई !!{element} नहीं है, तो $A = B$।
\end{prob}

\end{document}
